\chapter{Score Matching Based Detection}
\label{score_matching_detection}
\label{sec:score_matching_detection}
We now turn to the main methodological step: estimating the background score from target-free secondary data and plugging it into the LMP detector. Let
\begin{equation}
    w_1,\ldots,w_n \overset{\mathrm{i.i.d.}}{\sim} p_w
\end{equation}
denote the secondary samples.

To estimate the background score without evaluating the normalizing constant of $p_w$, we can use implicit score matching (ISM) introduced by Hyv\"arinen~\cite{hyvarinen2005estimation}:
\begin{equation}
    \min_\eta \frac{1}{n} \sum_i \left[
        \tfrac{1}{2}\|\score_\eta(w_i)\|^2
        + \operatorname{tr}\!\bigl(\nabla_w \score_\eta(w_i)\bigr)
    \right].
\end{equation}
This objective contains the divergence term, i.e. the trace of the Jacobian of $\score_\eta$, which is costly to evaluate and back-propagate through for neural models. We therefore turn to
\emph{denoising} score matching (DSM)~\cite{vincent2011connection},
which replaces that term by a Gaussian-perturbation regression with a fixed corruption
level $\sigma>0$. For each secondary sample we draw a corrupted copy
\begin{align}
    \tilde w_i &= w_i + \sigma\varepsilon_i,
    \nonumber\\
    \varepsilon_i&\sim \N(0,I_d),
\end{align}
and train $\score_\eta:\R^d\to\R^d$ by minimizing
\begin{equation}
    \hat\eta
    = \arg\min_\eta
    \frac1n\sum_{i=1}^n
    \left\|
    \score_\eta(\tilde w_i) - \frac{w_i-\tilde w_i}{\sigma^2}
    \right\|^2.
    \label{eq:DART_objective}
\end{equation}
The regression target $(w_i-\tilde w_i)/\sigma^2$ is the analytic score of the Gaussian
corruption kernel. At the population optimum,
\begin{align}
    \score_{\eta^*}(x)
    &= \nabla_x \log q_\sigma(x)\nonumber\\
    q_\sigma &= p_w * \mathcal{N}(0,\sigma^2 I_d),
\end{align}
so $\score_{\hat\eta}$ estimates the score of the noise-smoothed density $q_\sigma$, and
tends to $\nabla\log p_w$ as $\sigma\to0$ (Tweedie's formula~\cite{efron2011tweedie}%
). 

The noise level $\sigma$ governs a bias-variance tradeoff: a small $\sigma$ keeps $q_\sigma$ close to $p_w$, while a larger one stabilizes estimation and keeps the score informative away from the training samples.
Motivated by the analogy with kernel bandwidth selection~\cite{silverman1986density}, we set $\sigma$ proportional to the empirical RMS scale of the data, $\sigma^2 = \rho\,\hat{s}^2$ with $\hat{s}^2 = \operatorname{tr}(\widehat{\Sigma}_n)/d$, where $\widehat{\Sigma}_n$ is the empirical covariance of the background samples, computed in the normalized coordinates on which the score network operates.
The dimensionless factor $\rho$ is small, so $\sigma$ remains well below the data scale, and decreases as the number of secondary samples grows, just as a kernel bandwidth narrows with sample size.
After learning the score, we estimate the score covariance from the same secondary samples:
\begin{align}
    \widehat{\Psi}
    = \frac{1}{n}\sum_{i=1}^n
    \score_{\hat\eta}(w_i)\score_{\hat\eta}(w_i)^\top.
    \label{eq:psi_hat}
\end{align}
This yields the proposed DART detector
\begin{align}
    \widehat T_{\mathrm{DART}}(y)
    = \frac{-s^\top \score_{\hat\eta}(y)}
    {\sqrt{s^\top \widehat{\Psi}\, s}}.
    \label{eq:lmp_plugin}
\end{align}
While the ideal statistic is standardized under $\Hnull$, the null distribution of the plug-in statistic with a learned score is unknown, and its asymptotic normality cannot be established. We therefore set detection thresholds empirically, from an upper quantile of the statistic evaluated on the target-free secondary samples, $\{\widehat T_{\mathrm{DART}}(w_i)\}_{i=1}^n$.
The complete learning and test procedure is summarized in Fig.~\ref{fig:dart_pipeline}.

% Adapted from the MLSP 2026 poster (column 2): the DART pipeline at
% learning time and at test time.
\begin{figure}[htbp]
\centering
\resizebox{\textwidth}{!}{%
\begin{tikzpicture}[x=0.5cm, y=0.5cm, line cap=round, font=\footnotesize\linespread{1}\selectfont,
  dbox/.style={draw=pmuted!60, fill=white, rounded corners=1mm, line width=0.5pt,
               align=center, inner sep=1.3mm},
  darr/.style={-stealth, line width=0.7pt, pnavy!75}]
  % ================= LEARNING TIME =================
  \draw[pnavy!70, dashed, line width=0.6pt, rounded corners=1.5mm] (0,6.3) rectangle (27.2,11.4);
  \node[anchor=north west, font=\footnotesize\bfseries, text=pnavy] at (0.4,11.2) {Learning time};
  \node[dbox, anchor=west] (n1) at (2.3,8.3) {target-free\\ $w_1,\ldots,w_n \overset{\text{iid}}{\sim} p_w$};
  \node[dbox, right=0.6cm of n1] (n2) {score net $\psi_\eta$\\ DSM training};
  \node[dbox, right=0.6cm of n2] (n3) {statistics\\ $\widehat{T}_{\mathrm{DART}}(w_i)$};
  \node[dbox, right=0.6cm of n3] (n4) {threshold\\ $\gamma = \hat{q}_{\,1-\alpha}$};
  \draw[darr] (n1) -- (n2);
  \draw[darr] (n2) -- (n3);
  \draw[darr] (n3) -- (n4);
  % ================= TEST TIME =================
  \draw[pbrown!75, dashed, line width=0.6pt, rounded corners=1.5mm] (0,-0.1) rectangle (27.2,4.9);
  \node[anchor=north west, font=\footnotesize\bfseries, text=pbrown] at (0.4,4.7) {Test time};
  \node[dbox] (m4) at (n4.center |- 0,2.1) {declare $\Halt$ if\\ $\widehat{T}_{\mathrm{DART}}(y) \gtrless \gamma$};
  \node[dbox, left=0.6cm of m4] (m3) {statistic\\ $\widehat{T}_{\mathrm{DART}}(y)$};
  \node[dbox, left=0.6cm of m3] (m2) {score net $\psi_{\hat\eta}$\\ (frozen)};
  \node[dbox, left=0.6cm of m2] (m1) {test pixel\\ $y_i$};
  \draw[darr] (m1) -- (m2);
  \draw[darr] (m2) -- (m3);
  \draw[darr] (m3) -- (m4);
  % gamma feeds the test-time decision
  \draw[darr] (n4.south) -- (m4.north);
  \node[text=pnavy, anchor=west] at ([shift={(0.3,-1.2)}]n4.south) {$\gamma$};
\end{tikzpicture}}
\caption{The DART pipeline. At learning time the score network is trained
on the target-free secondary samples with the DSM objective
\eqref{eq:DART_objective}, the statistic $\widehat{T}_{\mathrm{DART}}$ is
evaluated on the same samples, and the detection threshold $\gamma$ is set
to its empirical upper quantile. At test time the frozen network is applied
to each pixel and the statistic \eqref{eq:lmp_plugin} is compared with
$\gamma$.}
\label{fig:dart_pipeline}
\end{figure}


\section{Special case: Gaussian background}
\label{sec:special}

The Gaussian case makes the role of score matching explicit. Suppose
$
w \sim \mathcal{N}(\mu,\Sigma)
$
and restrict the score model to the affine class
$
\psi_\eta(w)=A(w-\mu).
$
Since
\begin{equation}
\psi(w)=-\Sigma^{-1}(w-\mu),
\end{equation}
the affine class is correctly specified.

\noindent\textbf{ISM.}
Minimizing the implicit score-matching objective over the affine class recovers
the empirical Gaussian score ~\cite{hyvarinen2005estimation}:
\begin{equation}
\hat\psi_{\mathrm{ISM}}(y)=-\hat\Sigma_n^{-1}(y-\hat\mu_n),
\end{equation}
where $\hat\mu_n$ and $\hat\Sigma_n$ are the sample mean and covariance.
Substituting this into the plug-in LMP statistic yields the usual Gaussian
matched-filter form, i.e. the adaptive matched filter (AMF) ~\cite{robey1992cfar,kelly1986adaptive}.

\noindent\textbf{DSM.}
Denoising score matching estimates the score of the smoothed density
$
q_\sigma = p_w * \mathcal{N}(0,\sigma^2 I)
$
rather than that of $p_w$ itself \cite{vincent2011connection}. Equivalently,
\begin{equation}
\nabla_{\tilde w}\log q_\sigma(\tilde w)
=
\frac{\mathbb{E}[w\mid \tilde w]-\tilde w}{\sigma^2},
\label{eq:tweedie}
\end{equation}
which expresses the score in terms of the posterior mean under Gaussian
corruption \cite{efron2011tweedie}. For Gaussian clutter,
\begin{equation}
    q_\sigma=\mathcal{N}(\mu,\Sigma+\sigma^2 I)
\end{equation}
hence
\begin{equation}
\hat\psi_{\mathrm{DSM}}(y)
=
-(\hat\Sigma_n+R_n)^{-1}(y-\hat\mu_n+r_n).
\label{eq:gauss_DART_solution0}
\end{equation}
where $R_n$ and $r_n$ are empirical moments of the noise. Asymptotically, these residual moments vanish and the DSM score converges to
\begin{equation}
\hat\psi_{\mathrm{DSM}}(y)
\rightarrow
-(\hat\Sigma_n+\sigma^2 I)^{-1}(y-\hat\mu_n).
\label{eq:gauss_DART_solution}
\end{equation}
Thus, in the affine Gaussian setting, DSM replaces the inverse sample
covariance by a diagonally loaded inverse covariance. In adaptive detection,
this is the classical regularization used to stabilize matched-filter
statistics when $\hat\Sigma_n$ is ill-conditioned or the number of secondary
samples is small ~\cite{carlson1988covariance}. This connection explains the improved
small-sample robustness of linear DART in our experiments.
